% This is the tex file for the thesis/dissertation template

% =================================================================================================
% This is the preamble formatting for your thesis.
% It is called and used when using R Markdown and bookdown and knitting to PDF (as opposed to HTML
% or Word.
% =================================================================================================

% --- Main Formatting -----------------------------------------------------------------------------
% ----- Fonts and Paragraph Styling -----------------------------------------------------
\usepackage{fontspec}                             % Allows the use of system fonts. WARNING: This 
%                                                 % may require compiling with XeLaTex or LuaLaTex
\setmainfont{Times New Roman}                     % Set font to Times New Roman

\usepackage{setspace}                             % Provides commands for controlling line spacing
\onehalfspacing                                   % Set line spacing to 1.5

\usepackage{ragged2e}                             % Provides improved text alignment commands
\justifying                                       % Justify body text

\setlength{\parindent}{1.5em}                     % Indent the beginning of every paragraph, except 
%                                                 % the first

% ----- Text Emphasis -------------------------------------------------------------------
\usepackage[normalem]{ulem}                       % Provides underlining and strikeout commands. 
%                                                 % The command `\uline' is used in the Committee
%                                                 % Form Rmd and comes from this package.

% --- Header styling ------------------------------------------------------------------------------
\usepackage{titlesec}                             % Controls the style and spacing of section 
%                                                 % headers

% ----- Section header 1 styling --------------------------------------------------------
\titleformat{\section}[block]
  {%
    \normalfont
    \fontsize{14}{14}\selectfont                  % 14 pt font with 14 pt line height
    \bfseries                                     % Bold
    \centering%                                   % Center
  }
  {}                                              % Do not add section numbers before header title
  {0pt}                                           % No indentation to the left of the header title
  {}

\titlespacing*{\section}{0pt}{0pt}{0pt}           % No left indentation, no vertical space above 
%                                                 % the header, and no vertical space below the 
%                                                 % header

% ----- Section header 2 styling --------------------------------------------------------
\titleformat{\subsection}[block]
  {%
    \normalfont
    \fontsize{12}{12}\selectfont                  % 12 pt font with 12 pt line height
    \bfseries%                                    % Bold
  }
  {\thesubsection}                                % Include the subsection number
  {1em}                                           % Add 1em space between the subsection number 
  %                                               % and header title
  {}

\titlespacing*{\subsection}{0pt}{12pt}{0pt}       % No left indentation, 12 pt vertical space above
%                                                 % the header, and no vertical space below the 
%                                                 % header

% ----- Section header 3 styling --------------------------------------------------------
\titleformat{\subsubsection}[block]
  {%
    \normalfont
    \fontsize{12}{12}\selectfont                  % 12 pt font with 12 pt line height
    \bfseries%                                    % Bold
  }
  {\thesubsubsection}                             % Include the (level 2) subsection number
  {1em}                                           % Add 1em space between the subsection number 
  %                                               % and header title
  {}

\titlespacing*{\subsubsection}{0pt}{12pt}{0pt}    % No left indentation, 12 pt vertical space above
%                                                 % the header, and no vertical space below the 
%                                                 % header

% --- Tables of Contents --------------------------------------------------------------------------
\usepackage{tocloft}                              % Customizes the Table of Contents, list of 
%                                                 % figures, and list of tables

% ----- Add Leader Lines ----------------------------------------------------------------
\renewcommand{\cftsecleader}{%                    % Add dotted leader lines for level 1 headings in
  \cftdotfill{\cftdotsep}%                        % the Table of Contents
}
\renewcommand{\cftsubsecleader}{%                 % Add dotted leader lines for level 2 headings in
  \cftdotfill{\cftdotsep}%                        % the Table of Contents
}
\renewcommand{\cftsubsubsecleader}{%              % Add dotted leader lines for level 3 headings in
  \cftdotfill{\cftdotsep}%                        % the Table of Contents
}

% ----- Remove Titles -------------------------------------------------------------------
% Title are already added within the Rmd files, so the default titles are not needed
\renewcommand{\contentsname}{}                    % Remove the "Contents" title
\renewcommand{\listfigurename}{}                  % Remove the "List of Figures" title
\renewcommand{\listtablename}{}                   % Remove the "List of Tables" title

% ----- Add "Figure" in List of Figures -------------------------------------------------
\renewcommand{\cftfigpresnum}{Figure~}            % Add the word "Figure" before the figure numbers
%                                                 % in the List of Figures
\renewcommand{\cftfigaftersnum}{.}                % Add a period after the figure numbers in the 
%                                                 % List of Figures
\setlength{\cftfignumwidth}{5em}                  % Set the horizontal space reserved for figure 
%                                                 % numbers in the List of Figures. WARNING: This 
%                                                 % may need to be increased if the figure numbers
%                                                 % overlap with the caption text.

% ----- Add "Table" in List of Tables ---------------------------------------------------
\renewcommand{\cfttabpresnum}{Table~}             % Add the word "Table" before the table numbers 
%                                                 % in the List of Tables
\renewcommand{\cfttabaftersnum}{.}                % Add a period after the table numbers in the 
%                                                 % List of Tables
\setlength{\cfttabnumwidth}{5em}                  % Set the horizontal space reserved for table 
%                                                 % numbers in the List of Tables. WARNING: This 
%                                                 % may need to be increased if the table numbers
%                                                 % overlap with the caption text.

% --- Caption Styling -----------------------------------------------------------------------------
% ----- Font ----------------------------------------------------------------------------
\usepackage{caption}                              % Allows for customization of figure and table 
%                                                 % captions
\captionsetup{
  labelfont={bf,it},                              % Bold and italicize the labels
  font=it,                                        % Italicize the caption text 
  justification=justified,                        % Justify the caption
  singlelinecheck=false}                          % Do not justify single lines (if the caption is
%                                                 % just one line, then it won't force the 
%                                                 % justification alignment

% ----- Numbers -------------------------------------------------------------------------
\usepackage{amsmath}
\numberwithin{figure}{section}                    % Numbers figures within each level 1 section. 
%                                                 % For example, the first figure within Chapter 2 
%                                                 % will be numbered Figure 2.1

\numberwithin{table}{section}                     % Numbers tables within each level 1 section. For 
%                                                 % example, the first table within Chapter 2 will
%                                                 % be numbered Table 2.1

% --- Page Margins --------------------------------------------------------------------------------
\usepackage{geometry}                             % Controls the page size and margins
\geometry{
  letterpaper,                                    % Use page size letter
  top=1in,                                        % Top margin is 1 inch
  bottom=1in,                                     % Bottom margin is 1 inch
  left=1.25in,                                    % Left margin is 1.25 inch for binding
  right=1in                                       % Right margin is 1 inch
}

% Bibliography ------------------------------------------------------------------------------------
% Force hanging indent inside CSL References list environment
\usepackage{etoolbox}                             % Provides tools for modifying existing LaTeX 
%                                                 % commands and environments. It is being used 
%                                                 % here because the CSLReferences environment is
%                                                 % being redefined

\makeatletter                                     % Allows use of LaTeX internal commands that 
%                                                 % contain "@". This is needed when redefining 
%                                                 % some environments or commands
\renewenvironment{CSLReferences}[2]               % Redefines the CSLReferences environment 
%                                                 % generated by Pandox/bookdown. The `[2]` means
%                                                 % this environment takes two arguments from 
%                                                 % Pandoc: #1 and #2. #1 is the hanging indent and
%                                                 % #2 is the entry spacing. Instead of using them
%                                                 % directly, they are manually set in the next 
%                                                 % lines
 {\begin{list}{}{
  \setlength{\itemindent}{-1.5em}                 % Create the hanging indent by moving the first 
%                                                 % line left
  \setlength{\leftmargin}{1.5em}                  % and the other to the right.
  \setlength{\parsep}{0pt}                        % Remove the extra space between paragraphs  
%                                                 % within a reference entry
  \setlength{\itemsep}{0.2\baselineskip}          % Reduce the vertical space between references
 }}
 {\end{list}}
\makeatother                                      % Restore the usual meaning of @ in command names
